\section{The frame offset and phenomenology}\label{sec:pheno}

The kernel has no dynamics, so numbers can only enter through identifications of structural quantities with measured ones.  All such
identifications in this framework depend on one number, the offset between the matter frame and the light frame.

\subsection{The frame offset}

Let the matter frame and the tachyon frame sit at rapidities \(+\eta\) and \(-\eta\) on either side of the light frame, and write
\begin{equation}\label{eq:beta}
\beta_{\rm sym}=\tanh\eta,\qquad \gamma_L=\cosh2\eta,\qquad \beta_{\rm sym}=\sqrt{\frac{\gamma_L-1}{\gamma_L+1}} ,
\end{equation}
so that \(\gamma_L\) is the Lorentz factor between the two opposite frames.  We set
\begin{equation}\label{eq:gammaL}
\gamma_L=\frac{\gamma_M}{\gamma_T},\qquad \gamma_T=\frac{\ln2}{\pi\sqrt3}=0.127384,\qquad \gamma_M=0.237533,
\end{equation}
which gives \(\gamma_L=1.86470\), \(\beta_{\rm sym}=0.549406\) and \(\eta=0.6175\).  Both values on the right are values of the Barbero--Immirzi
parameter of loop quantum gravity.  The offset is therefore an input; nothing in Sections~\ref{sec:kernel}--\ref{sec:generation} fixes it.
Two structural values that suggest themselves do not apply: the value \(1/\sqrt3\) from a \(\mathbb Z_3\)-symmetric spinor on \(S^3\) describes a pure
state read from the light pole, whose populations \((\tfrac23,\tfrac13)\) differ from the matter frame's \((\tfrac34,\tfrac14)\), and a condition
\(\gamma_L=1/\beta_{\rm sym}\) equates a matter-frame readout with a tachyon-frame offset, i.e.\ \(\cosh2\eta=e^\eta\).

\subsection{A counting reading}

The two values in \eqref{eq:gammaL} solve one equation.  In the U(1) treatment of the black-hole horizon~\cite{abck,domagala,meissner} the
parameter \(\gamma\) satisfies
\begin{equation}\label{eq:counting}
\sum_j g(j)\,e^{-2\pi\gamma\sqrt{j(j+1)}}=1 .
\end{equation}
Spin \(j=\tfrac12\) alone with \(g=2\) gives \(\gamma_T\) exactly; all spins with \(g=2\) give \(\gamma_M\)~\cite{meissner}; all spins with \(g=2j+1\) give
\(\gamma=0.274067\)~\cite{ghosh} and would give \(\beta=0.6045\).  Under this reading \(\beta_{\rm sym}\) is the ratio of two counting rules: the
tachyon frame counts minimally, the matter frame counts all spins with weight two.  A single principle produces both rules.

\begin{proposition}[Uncertainty saturation]\label{prop:robertson}
In \(|j,m\rangle\) (\(\hbar=1\)), \(\Delta J_x\,\Delta J_y=\tfrac12\bigl(j(j+1)-m^2\bigr)\), and the Robertson bound \(\Delta J_x\,\Delta J_y\ge\tfrac12|m|\)
holds with equality exactly for \(m=\pm j\).  At \(j=\tfrac12\) every state saturates.
\end{proposition}

\begin{proof}
\(\langle J_x\rangle=\langle J_y\rangle=0\) and \(\langle J_x^2\rangle=\langle J_y^2\rangle=\tfrac12(j(j+1)-m^2)\).  The excess over the bound is
\(\tfrac12(j-|m|)(j+|m|+1)\).
\end{proof}

``Count only states that saturate the uncertainty bound'' thus gives the minimal counting for the tachyon frame and, at every spin, the
two states \(m=\pm j\) for the matter frame, which are the states counted in~\cite{meissner}.  The reading has a clear weakness: the
SU(2)-invariant treatment of the horizon~\cite{enp} is regarded as better founded than the U(1) treatment and leads to the value with
\(g=2j+1\).  The offset used here therefore rests on the less favoured counting, and the uncertainty principle is a reinterpretation of that
counting, not part of its derivation.

\subsection{Matter fraction}

The framework identifies the matter fraction with the square of the offset,
\begin{equation}\label{eq:omega}
\Omega_M=\beta_{\rm sym}^2=\frac{\gamma_M-\gamma_T}{\gamma_M+\gamma_T}=0.3018 ,
\end{equation}
to be compared with \(0.2975\pm0.0086\) from DESI DR2 baryon acoustic oscillations~\cite{desi} and \(0.3153\pm0.0073\) from Planck~\cite{planck}.  Since
\(\beta_{\rm sym}\) is constant and \(\Omega_M\) is not, \eqref{eq:omega} can hold at one epoch only; it is a statement about the present epoch,
and why now is open.  It concerns the share \(\Omega_\Lambda=1-\Omega_M\), not \(\Lambda\) itself: the frames differ in their counting, not in
\(G_N\), so \(\Lambda\) as a curvature is the same in all of them.  The comparison was made after the data were known.

\subsection{Masses}

The charged-lepton masses satisfy the Koide relation~\cite{koide}
\(\,(m_e+m_\mu+m_\tau)/(\sqrt{m_e}+\sqrt{m_\mu}+\sqrt{m_\tau})^2=\tfrac23\) to about \(10^{-5}\).  The framework does not derive it, but it gives its
form.  In the operator language of Section~\ref{sec:generation}, mass terms couple the matter frame to the tachyon frame and carry the
quantum numbers of a Higgs doublet.  A coupling with the full slot symmetry --- the cyclic shift \(\tau\) and the reflection \(\rho\) --- gives at
most two distinct masses.  Breaking \(\rho\) and keeping \(\tau\) gives a circulant coupling \(a+bP+\bar bP^2\) with eigenvalues
\begin{equation}\label{eq:brannen}
\sqrt{m_k}=a+2|b|\cos\Bigl(\theta+\frac{2\pi k}{3}\Bigr),\qquad k=0,1,2,
\end{equation}
and the Koide value corresponds to \(|b|=a/\sqrt2\).  Mixing requires breaking \(\tau\) itself.  Each observed property thus has its breaking
stage --- mass the frame coupling, splitting the breaking of \(\rho\), mixing the breaking of \(\tau\) --- while \(|b|/a\) and \(\theta\) remain
unexplained.

\subsection{Lepton mixing}

Let each particle kind read its mass coupling with a subgroup of \(A_4\supset\mathbb Z_3=\langle\tau\rangle\): Dirac couplings are circulant and
read with \(\mathbb Z_3\), Majorana neutrinos admit a further coupling and read with a \(\mathbb Z_2\).  This reproduces the known trimaximal pattern
TM2, \(\sin^2\theta_{12}=1/(3\cos^2\theta_{13})\approx0.341\), which lies \(3.6\sigma\) above the JUNO value~\cite{juno}.  A correction
\(D=\exp(\varepsilon W)\) with a fixed diagonal \(W\) lowers \(\theta_{12}\) into the measured range while \(\theta_{13}\) and \(\theta_{23}\) stay inside theirs.
Since \(D\) is boost-like, its natural parameter is a rapidity; with \(\varepsilon=\eta\) the model meets the data, with the value from
\(1/\sqrt3\) it does not, and in all cases that fit it gives \(\sin^2\theta_{23}\in[0.491,0.498]\).  The model rests on the \(A_4\) literature and has
few settings, and the choice \(\varepsilon=\eta\) was made after the target windows were known.  Read in reverse, the lepton data fix
\(\eta\in[0.611,0.647]\), i.e.\ \(\beta_{\rm sym}\in[0.545,0.570]\): the input value lies inside, the values from \(1/\sqrt3\), from \(\gamma_L=1/\beta_{\rm sym}\)
and from the SU(2) counting lie outside.

\subsection{Predictions}

The following predictions were written down with their test data and decision criteria and committed to a public repository before the
relevant data appeared~\cite{repo}.  Data known at the time of registration do not count.

\begin{center}
\small
\renewcommand{\arraystretch}{1.3}
\begin{tabular}{@{}>{\raggedright\arraybackslash}p{0.33\textwidth}>{\raggedright\arraybackslash}p{0.13\textwidth}>{\raggedright\arraybackslash}p{0.24\textwidth}>{\raggedright\arraybackslash}p{0.21\textwidth}@{}}
\toprule
prediction & registered & test data & refuted if\\
\midrule
\(\sin^2\theta_{12}=\beta_{\rm sym}^2/\cos^2\theta_{13}=0.3086\pm0.0002\) & 20 Oct 2026 & next JUNO result with \(\sigma\le0.004\) & \(|{\rm pull}|\ge3\)\\
\(\cot\theta_W=\gamma_L\) (on shell): \(m_W=m_Z\gamma_L/\sqrt{1+\gamma_L^2}=80.3615\pm0.0018\) GeV & 21 Oct 2026 & first world average after registration & \(|{\rm pull}|\ge3\)\\
\(\sin^2\theta_{23}\le0.54\) (normal ordering) & 24 Oct 2026 & first global fit or DUNE/Hyper-K result after registration & \(2\sigma\) interval wholly above \(0.54\)\\
\(\sin^2\theta_{23}\in[0.49,0.50]\) (with \(\varepsilon=\eta\)) & 24 Oct 2026 & as above & \(2\sigma\) interval misses \([0.49,0.50]\)\\
\bottomrule
\end{tabular}
\end{center}

The first prediction was registered with the rounded input \(\gamma_M=0.2375\) (\(\beta_{\rm sym}=0.549348\)); with the unrounded value it
moves to \(0.3087\), well inside its error criterion.  The first two relations were noticed during other computations and lie close to known values; the third was noticed after an evaluation;
the fourth rests on a reading chosen after the windows were known, although the value of \(\theta_{23}\) was not known when it was chosen.
Present global fits do not decide the last two: NuFIT gives \(\sin^2\theta_{23}=0.561\) without and \(0.470\) with Super-Kamiokande atmospheric
data~\cite{nufit}.  A confirmation of the fourth would test the mixing model as a whole, since other models also predict \(\theta_{23}\) near
\(45^\circ\).  The point of the table is less the chance of a hit than that the numbers were fixed in advance.
